\documentclass{article}
\PassOptionsToPackage{numbers}{natbib}
\usepackage[final,main]{neurips_2025}
\usepackage{amsmath,amsfonts,amssymb,amsthm}
\usepackage[utf8]{inputenc} % allow utf-8 input
\usepackage[T1]{fontenc}    % use 8-bit T1 fonts
\usepackage[hidelinks]{hyperref}       % hyperlinks
\usepackage{url}            % simple URL typesetting
\usepackage{booktabs}       % professional-quality tables
\usepackage{nicefrac}       % compact symbols for 1/2, etc.
\usepackage{microtype}      % microtypography
\usepackage{xcolor}         % colors
\usepackage{cleveref}
\newtheorem{theorem}{Theorem}[section]
\newtheorem{lemma}[theorem]{Lemma}
\newtheorem{proposition}[theorem]{Proposition}
\theoremstyle{definition}
\newtheorem{definition}[theorem]{Definition}
\newtheorem{example}[theorem]{Example}

\title{CS 683 AI Poker Agent Project Proposal}

\author{%
  Nathan Pearce \\
  \texttt{npearce@umass.edu} \And
  Albi Marini \\
  \texttt{albimarini@umass.edu} \AND
  Girirajan Soundar Rajan \\
  \texttt{gsoundarraja@umass.edu} \And
  Sumit Saurabh \\
  \texttt{ssaurabh@umass.edu}
}

\begin{document}
\maketitle

\section{Introduction and Project Plan}

Our project group will develop and implement the AI poker agent as described in the final project instructions. Poker is suitable for an AI application as it is a zero-sum game with hidden information, random events, and an exceedingly large search space, making brute-force search impractical. Our agent will combine depth-limited adversarial search with a heuristic evaluation function that estimates the value of non-terminal states using information such as hand strength, pot size, board context, and observed opponent behavior.

\subsection{Our Contributions}

To make the final project more novel, we will incorporate two research-inspired ideas beyond the basic fixed-abstraction poker bot. First, the agent will use dynamic abstraction, meaning that it will reason coarsely in routine situations but switch to fine-grained reasoning in strategically important situations, such as later betting rounds after the river with large pots. Second, the agent will use robust opponent modeling, where it tracks the opponent's betting patterns over repeated hands and cautiously adjusts play to exploit predictable behavior. These ideas will be evaluated by comparing the full agent against simpler baselines, such as a fixed-abstraction version and a version without opponent modeling, and we expect this comparison to show whether the added techniques improve decision quality/match performance.

\section{Role of Each Team Member}

We plan to divide the work evenly among each group member as follows:

Nathan will implement the final \texttt{declareaction()} pipeline and integrate all helper functions written by group members. Albi will code the helper functions for dynamic abstraction. Girirajan will code the helper functions for opponent modeling. Sumit will implement the offline self-play training for parameter tuning and agent comparisons. All four group members will jointly write the final report.

\end{document}